% 第3章第7节：本章小结。

\section{本章小结}
\label{sec:graph-optimization-summary}

本章的六个问题都在有限图上从满足结构要求的候选中挑选数值最好者，但选择的对象和数值的含义各不相同。最短路选一条路径，比较弧长之和；最小生成树选一组连接全部顶点的边，比较建造成本之和；最大流给每条弧分配流量，在容量上界内最大化源汇之间的净输送量；最小费用流在供需固定时比较输送费用；匹配选互不共享端点的边，只数配对个数；指派要求每个对象都完成配对，再比较配对成本。同样写在边旁的一个数，可以是长度、成本、容量或单位费用，它们不能互换解释：例如最小生成树上两点间的路径未必是最短路，最大流也不关心沿哪条路线输送更便宜。

各算法依靠的论证与成立条件也不同。有向无环图上的递推依靠拓扑序，前驱的距离在使用前已经确定；Bellman--Ford算法允许负边长，靠“至多$k$条弧”的递推反复修正标签，并能识别负环的影响；Dijkstra算法只在边长非负时成立，每次永久确定一个顶点。Kruskal算法与Prim算法依靠割性质，用交换论证说明每次选的边可以留在某棵最优树中。最大流的Ford--Fulkerson方法与Edmonds--Karp算法沿残量网络中的增广路径增加流量，前者在整数容量下有限终止，后者按最少弧数选路，对实容量也给出增广次数的界。网络单纯形法沿树上的基本环调整流量，靠词典序规则排除零步长换基造成的循环。增广匹配沿交替路径翻转选边状态；匈牙利算法在紧边上增广，在受阻时调整标号。

许多算法停止时都留下一份可以独立核对的证书，表~\ref{tab:graph-opt-summary}把它们与输入条件、主要开销并列。最短路的证书是逐弧距离不等式；最大流的证书是同值的流与割；最小费用流的证书是流量与节点势满足的逐弧符号条件；最大匹配的证书是单位容量网络中与匹配边数同值的割；指派的证书是成本等于标号之和的紧边指派。最小生成树的最优性则回到割性质：每条非树边都不轻于它在树上所对应路径中的边。证书的价值在于检查者不必相信求解过程，只需核对可行性和若干不等式；浮点计算中这些核对都要带容差，近似满足不能直接当作精确证书。

\begin{table*}[!htb]
  \centering\small
  \caption{本章六个问题的对比（$n$为顶点数，$m$为无向边数，$a$为有向弧数，$\nu$为二部图较小一侧的顶点数，$N$为指派方阵的阶数）。开销均为本章所述实现的运算量。}
  \label{tab:graph-opt-summary}
  \begin{tabularx}{\linewidth}{@{}l>{\raggedright\arraybackslash}X>{\raggedright\arraybackslash}X>{\raggedright\arraybackslash}X>{\raggedright\arraybackslash}p{30mm}@{}}
    \toprule
    问题 & 选择对象与目标 & 方法与成立条件 & 结果证书 & 主要开销\\
    \midrule
    最短路 & 一条路径，弧长之和最小 & DAG递推（无有向环）；Bellman--Ford（允许负边，检测负环）；Dijkstra（边长非负） & 逐弧距离不等式与前驱路线 & $O(n+a)$；$O(na)$；$O((n+a)\log n)$\\
    最小生成树 & 连接全部顶点的边集，成本之和最小 & Kruskal、Prim，依靠割性质 & 无环、覆盖各分量；非树边不轻于树上对应路径的边 & $O(m\log m)$；$O(n^2+m)$或$O((n+m)\log n)$\\
    最大流 & 各弧流量，源汇净输送量最大 & 残量增广；整数容量下任意选路有限终止，按最少弧数选路对实容量成立 & 同值的流与割 & Edmonds--Karp为$O(na^2)$\\
    最小费用流 & 各弧流量，供需固定时费用最小 & 网络单纯形法：可行树基、基本环调整、两阶段初始化、词典序防循环 & 流量与节点势的逐弧符号条件 & 每轮$O(n+a)$，总轮数无统一多项式界\\
    二部最大匹配 & 互不共享端点的边，边数最多 & 交替路径增广，限于二部图 & 单位容量网络中同值的割 & $O(\nu(n+m))$\\
    指派 & 完美匹配，成本之和最小 & 匈牙利算法：可行标号与紧边匹配 & 成本等于标号之和 & $O(N^3)$\\
    \bottomrule
  \end{tabularx}
\end{table*}

本章的两种结构会在后面再次用到。把全部边权取负，或把选边次序改为从大到小，最小生成树的算法就求出最大权生成树；第~\ref{sec:graph-tree-decomposition-construction}节构造团树时，以极大团为顶点、团的交集大小为边权，调用的正是最大权生成树。最小割则在第~\ref{sec:graph-pairwise-binary-cut}节中用于最小化由一元项和成对项组成的二值函数：满足相应条件的目标可以写成割容量加常数，于是最大流算法给出精确的最小化配置。复用之前都要先说明问题怎样转成当前的网络与目标，不能只因为都用到图就共享最优性结论。

本章选择的始终是路径、树、流、匹配这类具有特定组合结构的对象，算法的正确性也来自这些结构本身的性质。下一章转而用集合函数评价一般的子集，候选对象可以是顶点、边、特征或传感器。那里将区分两种情形：一种是评价函数的一般性质（如收益递减）给出贪心选择的近似保证，另一种是特殊的表示（如可以写成图割）带来精确的求解。
